\newpage\thispagestyle{hl}
\hlbanner{Part III $\cdot$ Reduced-order modelling}{8}{Projection-based reduction: POD, Galerkin and affine sweeps}
\hlsection{The question}
\begin{hlquote}You need the tip deflection of a beam for a thousand combinations of two stiffnesses and a load position. Each finite element solve is quick, but a thousand of them are not, and the answers are smooth functions of the parameters.\end{hlquote}
\vspace{8pt}\begin{hlsum}\textbf{Projecting onto a few snapshot modes leaves a small system to solve for each parameter set.} For a beam of \cpdNFree{} unknowns, the worst error on \cpdNTest{} unseen parameter sets is \cpdGalerkinErrSix{} with 6 modes, \cpdGalerkinErrTwelve{} with 12 and \cpdGalerkinErrLast{} with 24. The chapter explains POD and Galerkin models with an affine parameter split, judges them only on parameters not used in training, and computes the break-even number of queries. The basis is only as good as its snapshots, and no certified bound is given here.\end{hlsum}
\hlsection{Key concepts}
\begin{description}\setlength\itemsep{5pt}\setlength\parsep{0pt}\renewcommand{\descriptionlabel}[1]{\hspace{\labelsep}\sffamily\bfseries\color{hlnavy}#1}
\item[Snapshots] full-order solutions at training parameters
\item[POD basis] singular vectors of the snapshot matrix; the singular values bound the training error
\item[Galerkin projection] residual orthogonal to the basis, giving a small dense system
\item[Affine parameter dependence] a sum of fixed matrices, so the projection is done once
\item[Offline and online] build the model once, then solve at the cost of a small system
\item[Break-even queries] the number of queries at which the offline cost is repaid
\end{description}
\hlsection{After this chapter you can}
\begin{itemize}\setlength\itemsep{3pt}
\item build a POD--Galerkin model with an affine parameter split
\item judge it on parameters it was not trained on
\item compute the break-even number of queries from its offline and online costs
\end{itemize}
\newpage\thispagestyle{hl}
\noindent{\sffamily\small\color{hlgrey}Chapter 8 \quad Projection-based reduction: POD, Galerkin and affine sweeps}\par
\hlfig{\hlimg{ch08_pod}}{Singular values of the snapshot matrix (left) and error on unseen parameters against the number of modes (right).}
\hlsection{Headline result}
\begin{hlheadline}For a beam model of \cpdNFree{} unknowns, the worst error on \cpdNTest{} unseen parameter sets is \cpdGalerkinErrSix{} with 6 modes, \cpdGalerkinErrTwelve{} with 12 and \cpdGalerkinErrLast{} with 24.\end{hlheadline}
\hlsection{From the code}
\begin{lstlisting}[style=hlpy]
pod = PodBasis().fit(S, n_modes=8)
aff = AffineDecomposition([K1, K2], lambda mu: [mu[0], mu[1]]).project(pod.V)
Fr = pod.V.T @ F
# online: one 8 x 8 solve per parameter
q = np.linalg.solve(aff.assemble_reduced((2.0, 1.0)), Fr)
u = pod.V @ q
\end{lstlisting}
\hlsection{A stated limit}
\begin{hllimit}The basis is only as good as the snapshots: the error on unseen parameters is the only guide away from the training set, and no certified bound is given in this chapter.\end{hllimit}
